% ============================================================
% CAPÍTULO 4: IMPLEMENTACIÓN
% ============================================================

\section{Diseño de la Interfaz (Frontend Design)}
% Describe las decisiones de diseño visual previas a la implementación:
% - Modo oscuro, paleta de colores
% - Distribución del layout (chat principal, panel lateral)
% - Herramientas usadas (Figma, etc.)
% TODO: Redactar


\section{Desarrollo del Frontend}
% Detalla la implementación técnica del frontend:
% - Framework/librería usada (React, etc.)
% - Estructura de componentes
% - Gestión del estado
% - Renderizado de código (syntax highlighting)
% - Responsive design
% TODO: Redactar
% TODO: Incluir fragmentos de código relevantes si procede


\section{Desarrollo del Backend y de la API}
% Detalla la implementación del backend:
% - Framework usado (Node.js, Express, etc.)
% - Estructura de la API REST (endpoints principales)
% - Integración con el SDK de la API de LLM
% - Gestión de la base de datos
% - Autenticación y seguridad
% TODO: Redactar
% TODO: Incluir tabla con endpoints de la API


\section{Integración del Tutor Socrático}
% Describe cómo se integra el system prompt con la API del LLM:
% - Flujo de una conversación completa
% - Gestión del contexto y memoria de la conversación
% - Manejo de errores y fallbacks
% TODO: Redactar


\section{Ejemplo de Diálogo}
% Muestra un caso de uso real del tutor socrático en acción:
% - Interacción usuario ↔ tutor
% - Análisis de cómo se cumplen los objetivos pedagógicos
% TODO: Redactar
% \begin{itemize}
%     \item \textbf{Usuario:} ``...''
%     \item \textbf{Tutor Socrático:} ``...''
% \end{itemize}


\section{Pruebas y Validación}
% Describe las pruebas realizadas:
% - Tests unitarios
% - Tests de integración
% - Pruebas con usuarios reales (si aplica)
% - Métricas y resultados
% TODO: Redactar
